\chapter{Automation and Integration}
\label{ch:automation}

\section{Armada orchestration}

Armada is the multi-jail, compose-style interface. It can initialize a
manifest, merge several files, build missing roots, and apply lifecycle actions
to selected services.

\begin{lstlisting}[language=bash]
blackship armada init
blackship armada -f base.toml -f production.toml config --show
blackship armada -f base.toml -f production.toml up --detach
blackship armada ps
\end{lstlisting}

Later files override earlier values. Keep secrets and host-local paths in a
protected final override rather than the shared base manifest. Validate the
merged view, not only each input in isolation.

\section{JSON clients}

A management application should begin with:

\begin{enumerate}
  \item \command{blackship --format json capabilities};
  \item a schema-version check;
  \item \command{check} against the selected manifest;
  \item read-only status calls; and
  \item a structured dry run before offering an apply action.
\end{enumerate}

Invoke Blackship with a fixed executable path and explicit working directory.
Capture stdout and stderr separately. Decode the envelope only from the stream
documented for the exit case. Never parse the human message to determine an
error category.

Use \command{--no-escalate} when the parent application owns authorization, or
\command{--non-interactive} in unattended jobs. An interactive privilege prompt
inside a GUI or CI pipe is both confusing and prone to deadlock.

\section{Rust library}

The library API supports ephemeral jails without the CLI, parser, MCP server,
or telemetry runtime:

\begin{lstlisting}[language=TOML]
[dependencies]
blackship = { version = "0.2", default-features = false }
\end{lstlisting}

\begin{lstlisting}[language=Rust]
use blackship::{ContainerSpec, Exec, Host, Limits, Network, Selector};
use std::time::Duration;

let host = Host::new("/var/db/ci/blackship")
    .with_zfs("zroot", "ci/blackship");

host.reap_orphans(&Selector::prefix("job-"))?;
host.ensure_release("15.1-RELEASE")?;

let job = host.create(
    &ContainerSpec::new("job-42", "15.1-RELEASE")
        .network(Network::Inherit)
        .limits(Limits::new().memory("4g").max_processes(512)),
)?;

let result = job.exec(
    Exec::shell("make test")
        .cwd("/work")
        .env("CI", "true")
        .timeout(Duration::from_secs(900)),
)?;

job.destroy()?;
\end{lstlisting}

\command{Host} is a clonable description of layout and storage, not a remote
connection. Containers own independent host descriptions and can be used from
several threads. Dropping a container attempts cleanup; call \command{destroy}
when errors must be observed or \command{keep} when the instance should outlive
the handle.

A non-zero executed command is represented in the execution result, while
failure to create, enter, supervise, or clean the jail is a typed
\command{blackship::Error}. Library calls require appropriate privilege and
never re-execute through an escalation helper.

\section{Cargo features}

\begin{longtable}{L{0.20\textwidth}L{0.16\textwidth}L{0.54\textwidth}}
\toprule
Feature & Default & Adds \\
\midrule
\command{cli} & yes & Binary and usage parser; implies telemetry. \\
\command{mcp} & yes & MCP host, schemas, tracing, and stdio server; implies CLI. \\
\command{telemetry} & through CLI & OTLP export through rama. \\
\command{inspect} & no & State-machine introspection. \\
\bottomrule
\end{longtable}

\section{MCP server}

\command{blackship mcp} speaks Model Context Protocol over stdin/stdout. It is
an authoring and validation surface, not a privileged lifecycle daemon. Its
read-only tools expose host facts and the real manifest/Jailfile parsers.
File-producing tools render and validate a draft, then require operator
elicitation before an atomic write.

The gate fails closed: refusal, timeout, an unsupported client, missing
confirmation, a nonexistent destination directory, or failed validation leaves
no file behind. Replacing a file requires an additional overwrite decision.

MCP resources include a generated CLI manual, KDL command tree, JSON schemas,
Jailfile reference, and reserved jail-parameter reference. Because these are
generated from live code, use them when exact current syntax matters.

\section{CI pattern}

A CI adapter can map one job to one ephemeral jail:

\begin{enumerate}
  \item reap orphaned jobs with a narrow prefix or label selector;
  \item ensure the requested release or committed image exists;
  \item create a uniquely named jail and bounded mounts/resources;
  \item copy or mount the workspace;
  \item execute steps with explicit timeouts;
  \item copy reports out; and
  \item destroy the jail in a finally/defer path.
\end{enumerate}

Timeout supervision kills the workload from the host side so a process inside
the jail cannot evade it. Use read-only mounts for shared source or caches when
the job should not mutate host data.

\section{Versioning expectations}

Pin crate and binary versions in automation. Treat the JSON
\setting{schema\_version} independently from the Blackship package version.
Feature-detect optional behavior with \command{capabilities}, and never assume a
host has ZFS, PF, RCTL, or an authorized escalator solely because another host
did.
